\documentclass[../main.tex]{subfiles}
\begin{document}

\chapter{Memory Consistency}
\lessoninfo{
  video={Lesson 21, videos 1--14},
  textbook={H\&P \cite{hennessy2017} §5.6; \cite{mosberger1993}},
  keywords={memory consistency, cache coherence, program order, consistency model, sequential consistency, strong consistency, load replay, relaxed consistency, MSYNC, data race, weak consistency, release consistency},
}

Lesson 19 made every read of a memory location return the value most
recently written to that location, and Lesson 20 used locks and barriers to
coordinate the threads that update shared data.  Neither addresses the order
in which accesses to \emph{different} locations take effect, and an
out-of-order processor changes that order as a matter of course.  This lesson
shows that such reordering can produce results that no execution of the
program in its own order could produce, defines sequential consistency, the
guarantee that rules these results out, and explains how an out-of-order
processor provides it.  It then turns to relaxed consistency, in which the
program marks the places where order matters, and to data races, which decide
whether a program can tell the difference.

\section{Coherence and consistency}
\label{sec:l21-coherence}

Cache coherence\index{cache coherence} (\cref{sec:l19-definition}) constrains the accesses to one
memory location: a read of the location, by any core, returns the value of
the most recent write to that same location, by any core.\source{Lesson 21,
videos 1--2; H\&P §5.6.}  This is what allows threads to share data at all.
Coherence, however, treats each location on its own.  It does not say whether
a core that writes location A and then reads location B performs the two
accesses in this order as far as the other cores can tell, or whether two
writes of one core, to A and to B, become visible to the other cores in the
order in which the core made them.

\begin{definition}[Memory consistency]
  The \term{memory consistency}[メモリ一貫性] of a shared-memory system is its
  guarantee about the order in which the memory accesses of all cores,
  including accesses to different locations, take effect.  It determines
  which combinations of values the reads of a parallel program can return.
\end{definition}

A write takes effect when it becomes visible to reads; a read takes effect at
the moment whose memory contents it returns.  A precise statement of the
orders that a system allows is a \term{consistency model}[一貫性モデル].  This
lesson develops the most important model, sequential consistency, and its
relaxed alternatives; \cref{sec:l21-models} surveys further models.

The two properties have different origins.  Coherence is needed because
caches keep copies of memory: a system without caches has a single copy of
every location and is coherent by construction.  Consistency does not depend
on caches.  Even with a single copy of every location, a processor that
performs its memory accesses in a different order than the program specifies
changes which values its reads return---and, as the next section shows,
current processors do exactly that.  A coherent system is therefore not
necessarily consistent.

\section{Why consistency matters}
\label{sec:l21-matters}

Whether the order of accesses to different locations matters is best seen in
a small program (\cref{ex:l21-flags}).\source{Lesson 21, videos 3--5; H\&P
§5.6.}

\begin{example}[Two flags]
  \label{ex:l21-flags}
  Two cores execute the same code.  On core 0, R4 holds the address of a
  flag X and R5 the address of a flag Y; on core 1 the roles of X and Y are
  exchanged.  Both flags are initially 0.
  \begin{lstlisting}[language={[x86masm]Assembler}, numbers=none]
      ADDI R1, R0, 1    ; R1 <- 1
      SW   R1, 0(R4)    ; own flag <- 1
      LW   R2, 0(R5)    ; R2 <- other core's flag
      BNEZ R2, skip     ; other core has announced
      ...               ; work for one core only
skip: ...
\end{lstlisting}
  Each core first announces with its own flag that it is about to do the
  work, and does the work only if the other core has not announced.
\end{example}

If each core performs its two memory accesses in program order, at most one
core does the work.  Suppose core 0 reads Y as 0.  Then its read of Y took
effect before core 1's store to Y, and therefore also before core 1's read of
X, which follows that store.  Core 0's store to X precedes its own read of Y,
so it precedes core 1's read of X as well, and that read returns 1.  (If both
stores take effect before both reads, neither core does the
work.)\margin{The two-flag program is the core of Dekker's and Peterson's
mutual-exclusion algorithms (Lesson 20), which use ordinary loads and
stores.}

An out-of-order processor does not perform the accesses in program order.  A
store writes memory only when it commits, whereas a load reads memory as soon
as it has computed its address and found no earlier store to the same address
in the load-store queue (\cref{sec:l09-lsq}).  The store and the load of
\cref{ex:l21-flags} access different addresses, so the load can read memory
long before the store commits.\margin{The distance is bounded by the reorder
buffer: 224 entries on a Skylake core (Lesson 8).  Agner Fog,
microarchitecture manual.}  If this happens on both cores
(\cref{fig:l21-reorder}), both loads read 0 and both cores do the work.

\begin{figure}[htbp]
  \centering
  % Two out-of-order cores run the two-flag program.  In program order each
% core stores 1 into its own flag and then loads the other core's flag.
% Both loads read memory before either store commits, so both read 0.
% Arrows: data read from memory (up/down towards a core) and stores written
% to memory at commit (towards the memory row).
\begin{tikzpicture}[x=0.95cm,y=1cm, font=\footnotesize\sffamily, >={Stealth[length=1.6mm]},
  row/.style={ln-muted, line width=0.6pt},
  ev/.style={circle, fill=black, inner sep=1.3pt},
  lab/.style={font=\scriptsize, align=center, inner sep=1pt},
  mem/.style={draw, fill=ln-box, minimum height=0.5cm, inner sep=1pt, font=\scriptsize}]
  % ---- row names
  \node[anchor=east] at (-0.15,1.5) {\iflangja{コア}{core}~0};
  \node[anchor=east] at (-0.15,0) {\iflangja{メモリ}{memory}};
  \node[anchor=east] at (-0.15,-1.5) {\iflangja{コア}{core}~1};
  \draw[row] (0,1.5) -- (10.2,1.5);
  \draw[row] (0,-1.5) -- (10.2,-1.5);
  % ---- memory contents over time
  \node[mem, minimum width=5.0*0.95cm, anchor=west] at (0,0) {X $=$ 0, Y $=$ 0};
  \node[mem, minimum width=2.2*0.95cm, anchor=west] at (5.0,0) {X $=$ 1, Y $=$ 0};
  \node[mem, minimum width=3.0*0.95cm, anchor=west] at (7.2,0) {X $=$ 1, Y $=$ 1};
  % ---- core 0: labels above its row
  \node[ev] at (0.7,1.5) {};
  \node[lab, anchor=south] at (0.7,1.62) {\iflangja{SW X\\実行}{SW X\\executes}};
  \node[ev, ln-caution] at (2.4,1.5) {};
  \node[lab, anchor=south, text=ln-caution] at (2.4,1.62) {\iflangja{LW Y\\0 を読む}{LW Y\\reads 0}};
  \draw[->, ln-caution] (2.4,0.27) -- (2.4,1.4);
  \node[ev] at (5.0,1.5) {};
  \node[lab, anchor=south] at (5.0,1.62) {\iflangja{SW X\\コミット}{SW X\\commits}};
  \draw[->] (5.0,1.4) -- (5.0,0.27);
  \node[ev, ln-caution] at (8.4,1.5) {};
  \node[lab, anchor=south, text=ln-caution] at (8.4,1.62) {\iflangja{LW Y コミット，\\作業を実行}{LW Y commits,\\does the work}};
  % ---- core 1: labels below its row
  \node[ev] at (1.3,-1.5) {};
  \node[lab, anchor=north] at (1.3,-1.62) {\iflangja{SW Y\\実行}{SW Y\\executes}};
  \node[ev, ln-caution] at (3.6,-1.5) {};
  \node[lab, anchor=north, text=ln-caution] at (3.6,-1.62) {\iflangja{LW X\\0 を読む}{LW X\\reads 0}};
  \draw[->, ln-caution] (3.6,-0.27) -- (3.6,-1.4);
  \node[ev] at (7.2,-1.5) {};
  \node[lab, anchor=north] at (7.2,-1.62) {\iflangja{SW Y\\コミット}{SW Y\\commits}};
  \draw[->] (7.2,-1.4) -- (7.2,-0.27);
  \node[ev, ln-caution] at (9.4,-1.5) {};
  \node[lab, anchor=north, text=ln-caution] at (9.4,-1.62) {\iflangja{LW X コミット，\\作業を実行}{LW X commits,\\does the work}};
  % ---- time axis
  \draw[->, ln-muted] (0,-2.45) -- (10.2,-2.45)
    node[below left, inner sep=2pt, font=\scriptsize] {\iflangja{時間}{time}};
\end{tikzpicture}

  \caption{\Cref{ex:l21-flags} on two out-of-order cores.  Each core's load
    reads memory before its own store commits; both loads read 0, and both
    cores do the work.}
  \label{fig:l21-reorder}
\end{figure}

Every read in this execution is coherent: when each load reads its flag, the
flag really holds 0, and the stores take effect at commit, in program
order.\caution{No cache holds a stale value here, so no coherence protocol
can prevent this outcome: the cause is reordering, not stale data.}
The result is nevertheless impossible in any execution in program order.  What
the out-of-order cores have changed is the order of accesses to different
locations: each core's read of the other flag took effect before its own
store.

\subsection{Reordered loads in correct programs}
\label{sec:l21-ready}

The two-flag program may look contrived, but the same kind of reordering
breaks common and otherwise correct patterns.  One is a \emph{data-ready
flag}: a writer thread stores new data and then sets a flag, and a reader
thread waits until the flag is set and then uses the data.  With R4 holding
the address of the flag and R5 that of the data, the writer executes
\begin{lstlisting}[language={[x86masm]Assembler}, numbers=none]
      SW   R6, 0(R5)    ; data  <- new value
      SW   R7, 0(R4)    ; ready <- R7 (= 1)
\end{lstlisting}
and the reader executes
\begin{lstlisting}[language={[x86masm]Assembler}, numbers=none]
poll: LW   R1, 0(R4)    ; R1 <- ready
      BEQZ R1, poll     ; not set: poll again
      LW   R2, 0(R5)    ; R2 <- data
\end{lstlisting}
In program order, the reader's load of the data follows the load of the flag
that found it set, which follows the writer's store to the flag, which in turn
follows the writer's store of the data: the reader always obtains the new
data.  Branch prediction (Lesson 4) removes this guarantee.  The reader's core
predicts that the loop exits and executes the load of the data speculatively,
possibly before the load of the flag in the same iteration.
\Cref{tab:l21-ready} traces such an execution in which the data
holds 3 and the writer stores 8.  The load of the flag executes after the
writer has set it, so the prediction turns out to be correct, and the load of
the data commits with the old value.

\begin{table}[htbp]
  \centering\small
  \caption{The data-ready flag on an out-of-order reader core (core 0) and a
    writer core (core 1).  Memory contents after each step.}
  \label{tab:l21-ready}
  \begin{tabular}{@{}cclcc@{}}
    \toprule
    Step & Core & Event & ready & data \\
    \midrule
    1 & 0 & \texttt{LW data} executes on the predicted path: R2 $=$ 3 & 0 & 3 \\
    2 & 1 & \texttt{SW data} commits                               & 0 & 8 \\
    3 & 1 & \texttt{SW ready} commits                              & 1 & 8 \\
    4 & 0 & \texttt{LW ready} executes: R1 $=$ 1                   & 1 & 8 \\
    5 & 0 & \texttt{BEQZ} resolves: not taken, as predicted        & 1 & 8 \\
    6 & 0 & the three instructions commit with R2 $=$ 3            & 1 & 8 \\
    \bottomrule
  \end{tabular}
\end{table}

The same happens when a thread creates another thread, waits until that
thread has exited, and then uses the results it produced.  The wait ends when
a load of the thread's status, which the operating system writes when the
thread exits, finds the thread finished; the load of the results, which
follows in program order, can execute earlier, before the other thread has
written them.\margin{Thread libraries hide this: POSIX requires
\texttt{pthread\_join} to order the writes of the exiting thread before the
accesses of the joining one.}

In both patterns all caches are coherent, and the program is correct whenever
its accesses take effect in program order.  To be able to write such programs
at all, programmers need a guarantee about the order of accesses to different
locations---a guarantee beyond coherence.

\begin{keypoint}
  Coherence orders the accesses to each location separately.  Out-of-order
  execution and branch prediction let loads take effect before earlier
  accesses to other locations, which can produce results that no execution in
  program order could produce.
\end{keypoint}

\section{Sequential consistency}
\label{sec:l21-sc}

The most intuitive guarantee requires memory to behave as if the cores took
turns.\source{Lesson 21, video 6; H\&P §5.6; \cite{lamport1979}.}

\begin{definition}[Sequential consistency]
  A shared-memory system provides
  \term{sequential consistency}[逐次一貫性] if the result of every execution
  is the same as if the memory accesses of each core were performed in program
  order and the accesses of different cores were interleaved in some arbitrary
  order.
\end{definition}

The definition, which goes back to Lamport \cite{lamport1979}, has two parts.
Within one core the order is fixed: the accesses appear in the order that the
program specifies.  Across cores any interleaving is allowed; which one an
execution corresponds to depends on the timing of the cores, so a correct
program must produce a correct result for every interleaving.  The words
\enquote{as if} matter: the hardware need not actually perform the accesses in
such an order, as long as no program can tell the difference.

\begin{example}[Interleavings of the two-flag program]
  \label{ex:l21-interleavings}
  In \cref{ex:l21-flags} core 0 stores X ($w_0$) and then reads Y ($r_0$);
  core 1 stores Y ($w_1$) and then reads X ($r_1$).  Keeping the order of each
  core leaves six interleavings (\cref{tab:l21-interleavings}).%
  \tip{Interleavings that keep both program orders: $\binom{m+n}{m}$ for $m$
  and $n$ accesses, here $\binom{4}{2}=6$.}  In every one
  of them at least one read returns 1, so a sequentially consistent system
  never lets both cores do the work: the execution of \cref{fig:l21-reorder}
  is not sequentially consistent.
\end{example}

\begin{table}[htbp]
  \centering\small
  \caption{The six interleavings of the accesses in
    \cref{ex:l21-interleavings} that keep each core's program order.}
  \label{tab:l21-interleavings}
  \begin{tabular}{@{}clccl@{}}
    \toprule
    & Order & $r_0$ (Y) & $r_1$ (X) & Does the work \\
    \midrule
    1 & $w_0\;r_0\;w_1\;r_1$ & 0 & 1 & core 0 \\
    2 & $w_0\;w_1\;r_0\;r_1$ & 1 & 1 & neither \\
    3 & $w_0\;w_1\;r_1\;r_0$ & 1 & 1 & neither \\
    4 & $w_1\;w_0\;r_0\;r_1$ & 1 & 1 & neither \\
    5 & $w_1\;w_0\;r_1\;r_0$ & 1 & 1 & neither \\
    6 & $w_1\;r_1\;w_0\;r_0$ & 1 & 0 & core 1 \\
    \bottomrule
  \end{tabular}
\end{table}

\needspace{8\baselineskip}
\begin{note}[Strong consistency]
  Sequential consistency requires only that \emph{some} interleaving explain
  the results; the interleaving need not agree with the real time at which the
  accesses occur, so a read may return an old value although another core had
  already completed a write to that location.  A system in which the order
  must also respect real time---an access that completes before another begins
  comes first, so that every read returns the value most recently
  written---provides \term{strong consistency}[強い一貫性], called atomic
  consistency by Mosberger \cite{mosberger1993}.
\end{note}

\section{Implementing sequential consistency}
\label{sec:l21-implementing}

\subsection{The simple implementation}

The simplest way to provide sequential consistency is to remove the
\enquote{as if} from its definition.\source{Lesson 21, videos 7--8; H\&P
§5.6.}  Each core performs its memory accesses one at a time and in program
order: it starts an access only when all its earlier accesses have been
completed.  Coherence makes each access visible to all cores as it takes
effect, so the order in which the accesses of all cores take effect is itself
an interleaving of the kind the definition requires.

This implementation is correct but very slow.  An out-of-order core can still
reorder the instructions that do not access memory, but it can have only one
memory access in progress at a time: a load cannot execute until
every earlier store has committed, and every instruction that uses the loaded
value waits with it.  Memory accesses are, however, exactly the slow
operations that out-of-order execution most needs to overlap.  A cache miss
takes many cycles, and in a multiprocessor an access may in addition have to
wait for the bus and the coherence protocol (Lesson 19).\margin{Serialized,
each latency is paid in full: L1 4 cycles, L2 12, L3 44, memory 60--100~ns
on a Skylake-class core (Lesson 14).}  Limited to one
memory access at a time, a core loses most of the benefit of out-of-order
execution.

\subsection{Reordering loads safely}

A better implementation lets the core execute loads out of order, detects the
cases in which this could violate sequential consistency, and repairs them.
Two observations make this possible.
\begin{itemize}
  \item \textbf{Only loads need attention.}  Stores write memory at commit, in
    program order (\cref{sec:l09-write}), so they already take effect in
    program order.  A load that reads the location of an earlier store of the
    same core obtains the store's value from the load-store queue (Lesson 9).
  \item \textbf{The ROB records the program order.}  Whatever the order of
    execution, the ROB holds the instructions in program order, and a load
    that has executed early can be regarded as taking effect when it commits.
    The value it read is correct for that point unless another core has
    written the location in between.
\end{itemize}
The core therefore watches the coherence traffic of the other cores\margin{Intel's counters call these
memory-ordering machine clears (\texttt{MACHINE\_CLEARS}); each costs a
pipeline refill, about 17 cycles on Skylake (Lesson 4).} for the
locations of all loads that have executed but not yet committed.  Under
write-invalidate coherence, a write by another core arrives as an
invalidation of the block that contains the location
(\cref{sec:l19-write-invalidate}).  If such a write is observed, the load may
have read a stale value: the load and all instructions after it are cancelled
and executed again, exactly as after a branch misprediction
(\cref{sec:l08-recovery}).  This is always
possible because the load has not
committed.  Once the load commits, it has taken effect in program order, and
its location no longer needs to be watched.

\begin{figure}[htbp]
  \centering
  % Sequential consistency with reordered loads.  Core 0 executes, in program
% order, I1: LW ready and I2: LW data; core 1 stores data (8) and then
% ready (1).  I2 executes early and reads the old data (3).  While I2 is
% uncommitted (shaded window) core 0 watches for writes to data by other
% cores; the invalidation caused by core 1's store makes it replay I2.
\begin{tikzpicture}[x=0.95cm,y=1cm, font=\footnotesize\sffamily, >={Stealth[length=1.6mm]},
  row/.style={ln-muted, line width=0.6pt},
  ev/.style={circle, fill=black, inner sep=1.3pt},
  lab/.style={font=\scriptsize, align=center, inner sep=1pt},
  win/.style={fill=ln-accent!20, draw=none}]
  % ---- row names
  \node[anchor=east] at (-0.15,1.4) {\iflangja{コア}{core}~0};
  \node[anchor=east] at (-0.15,0) {\iflangja{コア}{core}~1};
  % ---- monitoring windows of I2 (executed, not yet committed)
  \fill[win] (0.6,1.27) rectangle (2.6,1.53);
  \fill[win] (6.6,1.27) rectangle (9.8,1.53);
  \draw[row] (0,1.4) -- (10.5,1.4);
  \draw[row] (0,0) -- (10.5,0);
  % ---- core 0
  \node[ev] at (0.6,1.4) {};
  \node[lab, anchor=south] at (0.6,1.62) {\iflangja{I2: LW data\\3 を読む}{I2: LW data\\reads 3}};
  \node[ev, ln-caution] at (2.6,1.4) {};
  \node[lab, anchor=south, text=ln-caution] at (2.6,1.62) {\iflangja{I2 を\\再実行へ}{I2 must\\be replayed}};
  \node[ev] at (5.0,1.4) {};
  \node[lab, anchor=south] at (5.0,1.62) {\iflangja{I1: LW ready\\1 を読む}{I1: LW ready\\reads 1}};
  \node[ev] at (6.6,1.4) {};
  \node[lab, anchor=south] at (6.6,1.62) {\iflangja{I2: LW data\\8 を読む}{I2: LW data\\reads 8}};
  \node[ev] at (8.2,1.4) {};
  \node[lab, anchor=south] at (8.2,1.62) {\iflangja{I1\\コミット}{I1\\commits}};
  \node[ev] at (9.8,1.4) {};
  \node[lab, anchor=south] at (9.8,1.62) {\iflangja{I2\\コミット}{I2\\commits}};
  % ---- core 1
  \node[ev] at (2.6,0) {};
  \node[lab, anchor=north] at (2.6,-0.12) {\iflangja{SW data\\（8）}{SW data\\(8)}};
  \node[ev] at (4.1,0) {};
  \node[lab, anchor=north] at (4.1,-0.12) {\iflangja{SW ready\\（1）}{SW ready\\(1)}};
  % ---- invalidation seen by core 0
  \draw[->, ln-caution, dashed] (2.6,0.1) -- (2.6,1.28)
    node[pos=0.45, left, lab, text=ln-caution] {\iflangja{data の\\無効化}{invalidation\\of data}};
  % ---- time axis
  \draw[->, ln-muted] (0,-0.95) -- (10.5,-0.95)
    node[below left, inner sep=2pt, font=\scriptsize] {\iflangja{時間}{time}};
\end{tikzpicture}

  \caption{Replaying a reordered load.  In program order core 0 executes I1:
    \texttt{LW ready} and then I2: \texttt{LW data}.  While I2 has executed
    but not committed (shaded), core 1 writes the data; its invalidation forces
    I2 to be executed again.}
  \label{fig:l21-replay}
\end{figure}

\Cref{fig:l21-replay} applies this to the data-ready flag of
\cref{tab:l21-ready}.  Core 0 executes the load of the data, I2, before the
load of the flag, I1, and reads the old value 3.  Core 1 then stores 8 into
the data.  The invalidation reaches core 0 while I2 is uncommitted, so I2 is
executed again and now reads 8.  I1, executed after core 1 has set the flag,
reads 1, and both loads commit with the values they would have returned in
program order.

A replay is needed only when another core writes a location between the
execution and the commit of a load that reads it.  This is uncommon in most
code, but frequent in loops that poll a location another core writes, such as
the reader of the data-ready flag.  Because invalidations refer to whole
blocks, a write to another location in the same block also causes a replay;
this costs time but never produces a wrong result.\margin{Blocks are 64
bytes on x86 and on most Arm cores, so unrelated variables share one: false
sharing (Lesson 19) causes replays as well as coherence misses.}

\begin{keypoint}
  Sequential consistency does not forbid reordering memory accesses; it
  forbids reordering that a program could observe.  An out-of-order core may
  execute loads early if it replays every load whose location another core
  writes before the load commits.
\end{keypoint}

\section{Relaxed consistency}
\label{sec:l21-relaxed}

Sequential consistency constrains every pair of accesses that a core makes,
whether or not the program depends on their order.\source{Lesson 21, videos
9--11; H\&P §5.6.}  For two accesses that a core makes in program order, the
first to location A and the second to location B, there are four kinds of
ordering: WR A $\to$ WR B, WR A $\to$ RD B, RD A $\to$ WR B and RD A $\to$ RD
B.  Each requires the first access to take effect before the second, as
observed by all cores (\cref{tab:l21-orderings}).  Sequential consistency
obeys all four, always.\margin{x86 relaxes only WR $\to$ RD; Arm, POWER and
RISC-V relax all four.  Intel SDM vol.~3A, \enquote{Memory Ordering}.}

\begin{table}[htbp]
  \centering\small
  \caption{The four orderings between two accesses of one core, first to A and
    then to B, and what a core that violates them may do.}
  \label{tab:l21-orderings}
  \begin{tabular}{@{}lp{4.3cm}p{4.1cm}@{}}
    \toprule
    Ordering & Violation & Example \\
    \midrule
    WR A $\to$ WR B & new value of B visible before new value of A
      & flag set before the data is stored \\
    WR A $\to$ RD B & B read before the write to A is visible
      & two flags (\cref{ex:l21-flags}) \\
    RD A $\to$ WR B & new value of B visible before A is read
      & a reply visible before the request has been read \\
    RD A $\to$ RD B & B read before A is read
      & data-ready flag (\cref{tab:l21-ready}) \\
    \bottomrule
  \end{tabular}
\end{table}

\begin{definition}[Relaxed consistency]
  A system with \term{relaxed consistency}[緩和一貫性] does not always obey
  some of these orderings for ordinary memory accesses, and provides a way for
  programs to enforce the order where their correctness depends on it.
\end{definition}

The two orderings that end in a read, WR $\to$ RD and RD $\to$ RD, are the
natural candidates for relaxation.  If a load may take effect before every
earlier access of its core to other locations, the core of
\cref{sec:l21-implementing} needs neither to watch coherence traffic nor to
replay loads because of other cores' writes: every load simply executes when
it is ready, as in \cref{fig:l21-reorder} and \cref{tab:l21-ready}.
Relaxing RD $\to$ RD alone only shortens the window in which a load must be
watched, because a load must still take effect after every earlier store of
its core.

\subsection{Synchronizing accesses}

If even two reads may be reordered, a program needs some way to obtain a
particular order where it matters.  A system with relaxed consistency
therefore distinguishes two kinds of operations.  Ordinary loads and stores
may be reordered as the model allows.  Special operations, which the program
uses where order matters, are never reordered.  The lecture calls such an
instruction \texttt{MSYNC}\index{MSYNC}: every memory access that precedes an
\texttt{MSYNC} in program order takes effect before any access that follows it
(\cref{fig:l21-msync}).\margin{x86 has such an instruction,
\texttt{MFENCE}\index{MFENCE}: the accesses before it become globally visible
before any that follow.}  Between two \texttt{MSYNC} instructions, accesses
may still be reordered.\caution{\texttt{MSYNC} is a \emph{full} fence.  Real
ISAs also have partial ones---x86 \texttt{SFENCE}, Arm \texttt{DMB
ISHST}---that order only some of the four pairs.}

\begin{figure}[htbp]
  \centering
  % Relaxed consistency with MSYNC.  The accesses of one core, in program
% order, are divided into groups by MSYNC instructions.  Accesses within a
% group may take effect in a different order; no access takes effect before
% an access that precedes it across an MSYNC.
\begin{tikzpicture}[x=1cm,y=1cm, font=\footnotesize\sffamily, >={Stealth[length=1.6mm]},
  acc/.style={draw, fill=ln-box, minimum width=1.05cm, minimum height=0.55cm, inner sep=1pt, font=\footnotesize\ttfamily},
  fence/.style={fill=ln-accent, minimum width=0.14cm, minimum height=1.1cm, inner sep=0pt},
  lab/.style={font=\scriptsize, align=center, inner sep=1pt}]
  % ---- accesses in program order
  \node[acc] (a1) at (0.0,0) {LW A};
  \node[acc] (a2) at (1.25,0) {SW B};
  \node[acc] (a3) at (2.5,0) {LW C};
  \node[fence] (f1) at (3.45,0) {};
  \node[acc] (a4) at (4.4,0) {LW D};
  \node[acc] (a5) at (5.65,0) {SW E};
  \node[acc] (a6) at (6.9,0) {LW F};
  \node[fence] (f2) at (7.85,0) {};
  \node[acc] (a7) at (8.8,0) {SW G};
  \node[acc] (a8) at (10.05,0) {LW H};
  \node[lab, anchor=south, text=ln-accent] at (f1.north) {MSYNC};
  \node[lab, anchor=south, text=ln-accent] at (f2.north) {MSYNC};
  % ---- reorderings allowed within a group
  \draw[<->, ln-tip] (a1.north) to[bend left=40] (a3.north);
  \draw[<->, ln-tip] (a4.north) to[bend left=40] (a6.north);
  \draw[<->, ln-tip] (a7.north) to[bend left=40] (a8.north);
  % ---- a reordering across MSYNC is not allowed
  \draw[->, ln-caution, dashed] (a5.south) to[bend left=35]
    node[pos=0.5, fill=white, inner sep=0.5pt, font=\small\bfseries] {$\times$} (a2.south);
  \node[lab, text=ln-caution, anchor=north] at (3.45,-1.15) {\iflangja{MSYNC をまたぐ\\順序変更は不可}{no reordering\\across MSYNC}};
  \node[lab, text=ln-tip, anchor=south] at (1.25,0.8) {\iflangja{グループ内は順序変更可}{reordering allowed}};
  \node[lab, text=ln-tip, anchor=south] at (5.65,0.8) {\iflangja{グループ内は順序変更可}{reordering allowed}};
  % ---- program order
  \draw[->, ln-muted] (-0.5,-1.95) -- (10.6,-1.95)
    node[below left, inner sep=2pt, font=\scriptsize] {\iflangja{プログラム順序}{program order}};
\end{tikzpicture}

  \caption{\texttt{MSYNC} divides the accesses of a core, in program order,
    into groups.  Accesses to different locations within a group may take
    effect in a different order, as far as the model allows; no access takes
    effect before an access in an earlier group.}
  \label{fig:l21-msync}
\end{figure}

With one \texttt{MSYNC} on each side, the data-ready flag is correct under
relaxed consistency.  The writer and the reader execute
\begin{lstlisting}[language={[x86masm]Assembler}, numbers=none]
; writer
      SW   R6, 0(R5)    ; data  <- new value
      MSYNC             ; the data is visible
      SW   R7, 0(R4)    ; ready <- 1
; reader
poll: LW   R1, 0(R4)    ; R1 <- ready
      BEQZ R1, poll     ; not set: poll again
      MSYNC             ; the load of ready is done
      LW   R2, 0(R5)    ; R2 <- data
\end{lstlisting}
On the reader, \texttt{MSYNC} makes the load of the flag that found the flag
set take effect before the load of the data.  On the writer, it makes the
store of the data take effect before the store to the flag, which a
relaxation of WR $\to$ WR could otherwise make visible first.  Because the
reader's load of the flag returned 1, the flag, and before it the data, had
been stored when that load took effect, and the load of the data comes later
still.  If only the orderings that end in a read are relaxed, the writer's two
stores already take effect in program order and the reader's \texttt{MSYNC}
alone suffices; the writer's \texttt{MSYNC} is needed when WR $\to$ WR may be
relaxed as well.

The flag is a simple form of synchronization, and the placement of the two
\texttt{MSYNC}s follows a general rule: place \texttt{MSYNC} \emph{after} an
operation that acquires the right to proceed and \emph{before} an operation
that releases it.\tip{Remember the placement as two fences facing
\emph{into} the critical section: nothing may leave it at either end.}  A
thread that uses a lock (Lesson 20) executes
\texttt{MSYNC} right after \texttt{lock}, so that the accesses of its critical
section take effect after the acquisition and see all writes of earlier lock
holders, and right before \texttt{unlock}, so that all writes of its critical
section have taken effect when the next thread acquires the lock.  Such
instructions normally belong inside the lock and barrier routines of a
synchronization library rather than in every program.\margin{Arm folds the
fence into the access itself: \texttt{LDAR} (acquire) and \texttt{STLR}
(release) order one direction only; RISC-V has \texttt{.aq}/\texttt{.rl}
bits.}

\section{Data races and consistency}
\label{sec:l21-races}

Whether relaxed consistency can change the result of a program depends on
how the program accesses shared data.\source{Lesson 21, video 12; H\&P §5.6;
\cite{adve1990}.}

\begin{definition}[Data race]
  Two accesses to the same memory location made by different cores form a
  \term{data race}[データ競合] if at least one of them is a write---RD $\to$
  WR, WR $\to$ RD or WR $\to$ WR---and they are not ordered by
  synchronization, that is, no synchronization operation determines which of
  them takes effect first.
\end{definition}

\begin{figure}[htbp]
  \centering
  % Data race versus accesses ordered by synchronization.  Left: a load and a
% store of A on different cores with nothing between them; either may take
% effect first.  Right: a barrier after the load on core 0 and before the
% store on core 1 forces the load to take effect first.
\begin{tikzpicture}[x=1cm,y=1cm, font=\footnotesize\sffamily, >={Stealth[length=1.6mm]},
  acc/.style={draw, fill=ln-box, minimum width=1.5cm, minimum height=0.5cm, inner sep=1pt, font=\footnotesize\ttfamily},
  bar/.style={draw, fill=ln-accent!20, minimum width=1.5cm, minimum height=0.5cm, inner sep=1pt, font=\scriptsize\ttfamily},
  hd/.style={font=\scriptsize, text=ln-muted},
  ttl/.style={font=\footnotesize\sffamily\bfseries}]
  % ================= left: data race
  \node[hd] at (0,0.6) {\iflangja{コア}{core}~0};
  \node[hd] at (3,0.6) {\iflangja{コア}{core}~1};
  \node[acc] (l0) at (0,0) {LW A};
  \node[acc] (s1) at (3,-1.1) {SW A};
  \draw[<->, ln-caution, line width=0.8pt] (l0.east) to[bend left=25] (s1.north);
  \node[font=\scriptsize, align=center, text=ln-caution] at (0.75,-1.1)
    {\iflangja{どちらが先かは\\タイミング次第}{either order,\\depending on timing}};
  \node[ttl, text=ln-caution] at (1.5,-2.5) {\iflangja{データ競合}{data race}};
  % ================= right: ordered by a barrier
  \begin{scope}[xshift=6.2cm]
    \node[hd] at (0,0.6) {\iflangja{コア}{core}~0};
    \node[hd] at (3,0.6) {\iflangja{コア}{core}~1};
    \node[acc] (l0b) at (0,0) {LW A};
    \node[bar] (b0) at (0,-0.8) {BARRIER};
    \node[bar] (b1) at (3,-0.8) {BARRIER};
    \node[acc] (s1b) at (3,-1.6) {SW A};
    \draw[ln-accent, densely dashed] (b0.east) -- (b1.west);
    \draw[->, ln-tip, line width=0.8pt] (l0b.south east) -- (s1b.north west);
    \node[ttl, text=ln-tip] at (1.5,-2.5) {\iflangja{データ競合なし}{no data race}};
  \end{scope}
  % ---- time runs downwards
  \draw[->, ln-muted] (-1.1,0.3) -- (-1.1,-2.0)
    node[below, font=\scriptsize, inner sep=2pt] {\iflangja{時間}{time}};
\end{tikzpicture}

  \caption{Left: nothing orders the load and the store of A, a data race.
    Right: a barrier after the load and before the store makes the load take
    effect first.}
  \label{fig:l21-race}
\end{figure}

On the left of \cref{fig:l21-race}, a load of A on core 0 and a store to A on
core 1 have nothing between them; which of them takes effect first depends on
the timing of the cores, and the load returns the old or the new value
accordingly.  On the right, core 0 reaches a barrier (\cref{sec:l20-barrier})
after its load, and core 1 reaches the same barrier before its store.  No
thread leaves a barrier before all threads have arrived, so the load always
takes effect before the store, and there is no data race.  Accesses protected
by the same lock are ordered as well: one critical section ends with
\texttt{unlock} before the other begins with \texttt{lock}.\margin{Race
detectors (ThreadSanitizer, Helgrind) search for exactly this: two
conflicting accesses that no lock, barrier or fence orders.}

The importance of data races lies in the following property
\cite{adve1990}.
\begin{itemize}
  \item A \textbf{data-race-free program} produces the same results as under
    sequential consistency under every consistency model that recognizes its
    synchronization: one in which the synchronization operations are
    themselves sequentially consistent and each of them separates the earlier
    accesses of its core from the later ones, as \texttt{MSYNC} does.
    Reordering accesses to different locations, or reads of the same location,
    does not change any value that is read, and every pair of accesses whose
    order could change a value is ordered by synchronization, which such a
    model keeps in order.
  \item A program \textbf{with data races} has no such guarantee.  On a system
    with relaxed consistency anything can happen, including results that no
    interleaving explains.
\end{itemize}

A program can therefore be developed and debugged under sequential
consistency, where even a program with data races produces results that some
interleaving explains, and then run under relaxed consistency for speed once
synchronization has removed the races.\sidenote{Programming languages turned
this into a contract.  C++11 and Java guarantee that a data-race-free
program behaves as under sequential consistency, and in C++ a data race
makes the whole program undefined.  The synchronization these models
recognize is \texttt{std::atomic}---its default is
\texttt{memory\_order\_seq\_cst}, with release/acquire as the cheaper
option---and the locks built on it \cite{boehm2008}.}

The examples of this lesson all contain data races.  In
\cref{ex:l21-flags} each core reads a flag that the other core writes without
synchronization.  The data-ready flag of \cref{tab:l21-ready} is a form of
synchronization that the memory system does not recognize: under sequential
consistency the flag orders the reader's load of the data after the writer's
store, but a system with relaxed consistency does not know this.  The
\texttt{MSYNC}s of \cref{sec:l21-relaxed} make that ordering visible to it, so
the load of the data no longer races with the store; the polling loads of the
flag themselves remain unsynchronized, but they can only delay the reader,
because the flag changes only from 0 to 1.  The lost increment of
\cref{tab:l20-race} is also a data race.  Under sequential consistency its
possible results are at least limited to those of the interleavings of the two
threads; under relaxed consistency even that limit disappears.
\caution{A program that works on a sequentially consistent system can still
fail under relaxed consistency if it contains a data race.}

\begin{keypoint}
  A data-race-free program whose synchronization the memory system recognizes
  cannot tell relaxed consistency from sequential consistency, so relaxing
  consistency costs a properly synchronized program nothing.  A program with
  data races can observe anything.
\end{keypoint}

\section{Consistency models}
\label{sec:l21-models}

Sequential consistency and the relaxed consistency of \cref{sec:l21-relaxed}
are two points in a wide range of consistency models.\source{Lesson 21, videos
13--14; \cite{mosberger1993}.}  The models differ in which orderings they
relax, and in whether and how they distinguish synchronization from other
accesses.  Some models treat all accesses alike; from stronger to weaker:
\begin{description}
  \item[Strong consistency] Every read returns the value most recently written
    in real time (\cref{sec:l21-sc}).
  \item[Sequential consistency] Some interleaving that keeps the program order
    of each core explains all results.
  \item[Causal consistency] Under \term{causal consistency}[因果一貫性] all
    cores must see causally related writes---two writes of the same core, or a
    write made by a core that had read, directly or indirectly, the result of
    the earlier one---in the same order; concurrent writes may be seen in
    different orders by different cores.
  \item[Processor consistency]\index{processor consistency} The writes of one
    core are seen by all cores in the order in which that core made them, but
    different cores may see the writes of different cores in different orders,
    and a core's read may take effect before its own earlier write to another
    location.
\end{description}
Neither causal nor processor consistency forbids the execution of
\cref{fig:l21-reorder}: the two writes are concurrent, so each core may see
the other core's write only after its own read.

Other models distinguish synchronization, and a data-race-free program that
uses the synchronization they recognize behaves under them as under
sequential consistency:
\begin{description}
  \item[Weak consistency] Under \term{weak consistency}[弱い一貫性] accesses to
    synchronization variables, such as locks, are sequentially consistent and
    act like \texttt{MSYNC}: every earlier access takes effect before them and
    no later access does.  All other accesses may be reordered freely between
    them.
  \item[Release consistency]\index{release consistency} Synchronization is
    split into acquiring and releasing, as in \texttt{lock} and
    \texttt{unlock}.  An acquire keeps later accesses from taking effect before
    it, and a release waits until all earlier accesses have taken effect, so
    ordinary accesses are never reordered across the boundaries of a critical
    section---the placement of \texttt{MSYNC} in
    \cref{sec:l21-relaxed}.\margin{Release consistency was defined in 1990
    for the Stanford DASH multiprocessor, the machine that made the
    acquire/release split a hardware notion.}
  \item[Lazy release consistency]\index{lazy release consistency} Release
    consistency implemented lazily, originally for distributed shared memory
    built in software: the updates made in a critical section are transferred
    to another node only when that node accesses the data after acquiring the
    lock, not at every release.\sidenote{Lazy release consistency
    \cite{keleher1992} was designed for software distributed shared memory
    on a network of workstations, where propagating an update means sending
    a message: the TreadMarks system transfers a page only when a node that
    has acquired the lock touches it.}
  \item[Scope consistency]\index{scope consistency} Each lock defines a scope:
    a thread that enters a critical section sees the updates made in earlier
    critical sections of the same lock, but not necessarily those made under
    other locks \cite{iftode1996}.
\end{description}
% keeps this paragraph with the lesson summary, which alone would fill a page
\needspace{13\baselineskip}
\noindent The list is not complete, and new models continue to be proposed.  What they
share is the arrangement of this lesson: the system may reorder the accesses
that its model leaves unordered, and a program that needs a particular order
obtains it through synchronization.

\needspace{18\baselineskip}
\begin{keypoint}[Lesson summary]
  Coherence orders the accesses to one location; memory consistency concerns
  the order of accesses to different locations, which out-of-order execution
  and branch prediction change.  Sequential consistency requires every
  execution to look as if each core's accesses were performed in program
  order and those of different cores interleaved.  Performing accesses one at
  a time provides it but ruins out-of-order execution; instead, a core can
  execute loads early and replay any load whose location another core writes
  before the load commits.  Relaxed consistency does not always keep the
  orderings between accesses and gives programs synchronizing operations, such
  as \texttt{MSYNC}, placed after acquiring and before releasing.  A program
  without data races---conflicting accesses on different cores not ordered by
  synchronization---behaves as under sequential consistency in any model that
  recognizes its synchronization; a program with data races can observe
  anything.  Weak, release, lazy release and scope consistency differ in how
  they recognize synchronization, and causal and processor consistency in
  which orders of writes all cores must agree on.
\end{keypoint}

\end{document}
